\documentclass[ignorenonframetext,]{beamer}
\setbeamertemplate{caption}[numbered]
\setbeamertemplate{caption label separator}{: }
\setbeamercolor{caption name}{fg=normal text.fg}
\beamertemplatenavigationsymbolsempty
\usepackage{lmodern}
\usepackage{amssymb,amsmath}
\usepackage{ifxetex,ifluatex}
\usepackage{fixltx2e} % provides \textsubscript
\ifnum 0\ifxetex 1\fi\ifluatex 1\fi=0 % if pdftex
  \usepackage[T1]{fontenc}
  \usepackage[utf8]{inputenc}
\else % if luatex or xelatex
  \ifxetex
    \usepackage{mathspec}
  \else
    \usepackage{fontspec}
  \fi
  \defaultfontfeatures{Ligatures=TeX,Scale=MatchLowercase}
\fi
% use upquote if available, for straight quotes in verbatim environments
\IfFileExists{upquote.sty}{\usepackage{upquote}}{}
% use microtype if available
\IfFileExists{microtype.sty}{%
\usepackage{microtype}
\UseMicrotypeSet[protrusion]{basicmath} % disable protrusion for tt fonts
}{}
\newif\ifbibliography
\hypersetup{
            pdftitle={Lecture 10: GLMs: Poisson Regression, Overdispersion},
            pdfauthor={Nick Reich / Transcribed by Daveed Goldenberg, edited by Josh Nugent},
            pdfborder={0 0 0},
            breaklinks=true}
\urlstyle{same}  % don't use monospace font for urls

% Prevent slide breaks in the middle of a paragraph:
\widowpenalties 1 10000
\raggedbottom

\AtBeginPart{
  \let\insertpartnumber\relax
  \let\partname\relax
  \frame{\partpage}
}
\AtBeginSection{
  \ifbibliography
  \else
    \let\insertsectionnumber\relax
    \let\sectionname\relax
    \frame{\sectionpage}
  \fi
}
\AtBeginSubsection{
  \let\insertsubsectionnumber\relax
  \let\subsectionname\relax
  \frame{\subsectionpage}
}

\setlength{\parindent}{0pt}
\setlength{\parskip}{6pt plus 2pt minus 1pt}
\setlength{\emergencystretch}{3em}  % prevent overfull lines
\providecommand{\tightlist}{%
  \setlength{\itemsep}{0pt}\setlength{\parskip}{0pt}}
\setcounter{secnumdepth}{0}
%       ************************************************
%       **        LaTeX preamble to be used with all 
%	**        statsTeachR labs/handouts.
%
%	Author: Nicholas G Reich
%	Last modified: July 2017
%	************************************************

%\documentclass[table]{beamer}

%	Set theme (a nice plain one)
\usetheme{Malmoe}

%	Use named colors, set main color of theme
%		to match Web site color:
\definecolor{MainColor}{RGB}{10, 74, 109}
\colorlet{MainColorMedium}{MainColor!50}
\colorlet{MainColorLight}{MainColor!20}
\usecolortheme[named=MainColor]{structure} 

%	For tables
%[dvipsnames] [table]
\usepackage{xcolor}

%% calling tabu.sty, assuming a particular directory structure
\usepackage{../../slide-includes/tabu}	% Even fancier than tabulary
\usepackage{multirow}

%	Just for the degree symbol
\usepackage{textcomp}

%	Get rid of footline (page, author, etc. on each slide)
\setbeamertemplate{footline}{}
%	Get rid of navigation buttons
\setbeamertemplate{navigation symbols}{}

%	Make footnotes not ugly
\usepackage{hanging}
\setbeamertemplate{footnote}{\raggedright\hangpara{1em}{1}\makebox[1em][l]{\insertfootnotemark}\footnotesize\insertfootnotetext\par}

%	Text style for code snippets inline in text:
\newcommand{\codeInline}[1]{\texttt{#1}}

%	Text style for emphasis stronger than \emph:
%		(Note, this doesn't toggle the way \emph does.
%			(Note, this can be done, didn't seem worth the trouble.))
\newcommand{\strong}[1]{{\bfseries{#1}}}


%        ******	Define title page	**********************
\setbeamertemplate{title page}{
	{\color{MainColor}
	% There must be a better way than this -vspace at
	%	 the top and bottom of the page to reduce the 
	%	 bottom margin, but I can't find one that works.
	%\vspace{-6em}

% 	% Go to a lot of trouble to get the title in a
% 	%	nice box, since customizing a beamer block
% 	%	does not entirely work here (I don't know why)
	\newlength{\titleBoxWidth}
	\setlength{\titleBoxWidth}{\textwidth}
	\addtolength{\titleBoxWidth}{-2.0em}
	\setlength{\fboxsep}{1.0em}
	\setlength{\fboxrule}{0pt}
	\fcolorbox{MainColor!25}{MainColor!25}{
		\parbox{\titleBoxWidth}{
			\raggedright
			\LARGE\textbf{\inserttitle}
		}	% end parbox
	}	% end fcolorbox

	\vfill
	\small{Author: \insertauthor}
	\vspace{\baselineskip}

	\small{Course: \underline{\href{http://nickreich.github.io/cda}{Categorical Data Analysis}} (BIOSTATS 743)}

%	\small{\Instructor}
%	\vspace{\baselineskip}

%	\small{\emph{This material is part of the \strong{statsTeachR} project}}

	\vfill
	
	\tiny{\emph{Made available under the \underline{\href{http://creativecommons.org/licenses/by-sa/4.0/}{Creative Commons Attribution-ShareAlike 4.0 International License}.}} \hfill \includegraphics[height=1em]{../../slide-includes/by-sa-compact.png}
 }


		\vspace{-15em}

	}	% end color
	\clearpage
}	% end define title page

\newcommand{\myitem}{\item[\tiny$\blacksquare$]}

% bold faced lower-case letters
\def\ba{\boldsymbol{a}}
\def\bb{\boldsymbol{b}}
\def\bc{\boldsymbol{c}}
\def\bd{\boldsymbol{d}}
%        \def\be{\boldsymbol{e}}
\def\bdf{\boldsymbol{f}}
\def\bg{\boldsymbol{g}}
\def\bh{\boldsymbol{h}}
%	\def\bi{\boldsymbol{i}}
\def\bj{\boldsymbol{j}}
\def\bk{\boldsymbol{k}}
\def\bl{\boldsymbol{l}}
\def\bm{\boldsymbol{m}}
\def\bn{\boldsymbol{n}}
\def\bo{\boldsymbol{o}}
\def\bp{\boldsymbol{p}}
\def\bq{\boldsymbol{q}}
\def\br{\boldsymbol{r}}
\def\bs{\boldsymbol{s}}
\def\bt{\boldsymbol{t}}
\def\bu{\boldsymbol{u}}
\def\bv{\boldsymbol{v}}
\def\bw{\boldsymbol{w}}
\def\bx{\boldsymbol{x}}
\def\by{\boldsymbol{y}}
\def\bz{\boldsymbol{z}}

% bold faced upper-case letters
\def\bA{\boldsymbol{A}}
\def\bB{\boldsymbol{B}}
\def\bC{\boldsymbol{C}}
\def\bD{\boldsymbol{D}}
\def\bE{\boldsymbol{E}}
\def\bF{\boldsymbol{F}}
\def\bG{\boldsymbol{G}}
\def\bH{\boldsymbol{H}}
\def\bI{\boldsymbol{I}}
\def\bJ{\boldsymbol{J}}
\def\bK{\boldsymbol{K}}
\def\bL{\boldsymbol{L}}
\def\bM{\boldsymbol{M}}
\def\bN{\boldsymbol{N}}
\def\bO{\boldsymbol{O}}
\def\bP{\boldsymbol{P}}
\def\bQ{\boldsymbol{Q}}
\def\bR{\boldsymbol{R}}
\def\bS{\boldsymbol{S}}
\def\bT{\boldsymbol{T}}
\def\bU{\boldsymbol{U}}
\def\bV{\boldsymbol{V}}
\def\bW{\boldsymbol{W}}
\def\bX{\boldsymbol{X}}
\def\bY{\boldsymbol{Y}}
\def\bZ{\boldsymbol{Z}}

% bold faced greek letters
\def\balpha{\boldsymbol{\alpha}}
\def\bbeta{\boldsymbol{\beta}}
\def\bgamma{\boldsymbol{\gamma}}
\def\bdelta{\boldsymbol{\delta}}
\def\bepsilon{\boldsymbol{\epsilon}}
\def\bvarepsilon{\boldsymbol{\varepsilon}}
\def\bzeta{\boldsymbol{\zeta}}
\def\bdeta{\boldsymbol{\eta}}
\def\btheta{\boldsymbol{\theta}}
\def\biota{\boldsymbol{\iota}}
\def\bkappa{\boldsymbol{\kappa}}
\def\blambda{\boldsymbol{\lambda}}
\def\bmu{\boldsymbol{\mu}}
\def\bnu{\boldsymbol{\nu}}
\def\bxi{\boldsymbol{\xi}}
\def\bomicron{\boldsymbol{\omicron}}
\def\bpi{\boldsymbol{\pi}}
\def\brho{\boldsymbol{\rho}}
\def\bsigma{\boldsymbol{\sigma}}
\def\btau{\boldsymbol{\tau}}
\def\bupsilon{\boldsymbol{\upsilon}}
\def\bphi{\boldsymbol{\phi}}
\def\bchi{\boldsymbol{\chi}}
\def\bpsi{\boldsymbol{\psi}}
\def\bomega{\boldsymbol{\omega}}
\def\bSigma{\boldsymbol{\Sigma}}

% environments
\newcommand{\be}{\begin{enumerate}}
\newcommand{\ee}{\end{enumerate}}
\newcommand{\bi}{\begin{itemize}}
\newcommand{\ei}{\end{itemize}}
\newcommand{\beqa}{\begin{eqnarray*}}
\newcommand{\eeqa}{\end{eqnarray*}}
\newcommand{\beqn}{\begin{eqnarray}}
\newcommand{\eeqn}{\end{eqnarray}}

\newtheorem{thm}{Theorem}
\newtheorem{cor}[thm]{Corollary} 
\newtheorem{prop}[thm]{Proposition}
\newtheorem{res}[thm]{Result}

% distributions, traces, norms
\newcommand{\IG}[2]{\mbox{IG}\left[ #1, #2 \right] }
\newcommand{\N}[2]{\mbox{N}\left[ #1, #2 \right] }
\newcommand{\MVN}[2]{\mbox{MVN}\left[ #1, #2 \right] }
\newcommand{\tr}[1]{\mbox{tr}\left[ #1 \right] }
\newcommand{\norm}[1]{\left|\left|#1\right|\right|}

% expected value & variance
\newcommand{\ev}{\mbox{E}}
\newcommand{\var}{\mbox{Var}}
\newcommand{\cov}{\mbox{Cov}}

% notation for functional regression
\newcommand{\M}{\boldsymbol{M}_{\psi\phi}}


% notation for VB updates
\newcommand{\1}{{\mathbf{1}}}

\newcommand{\psub}{\underline{p}}
\newcommand{\evq}{\ev_{q^{*}}}

\newcommand{\Ay}{A_{\bY}}
\newcommand{\By}{B_{\bY}}
\newcommand{\Ax}{A_{\bX}}
\newcommand{\Bx}{B_{\bX}}
\newcommand{\Ab}{A_{\bb}}
\newcommand{\Bb}{B_{\bb}}
\newcommand{\Ag}{A_{\bg}}
\newcommand{\Bg}{B_{\bg}}
\newcommand{\Al}{A_{\lambda}}
\newcommand{\Bl}{B_{\lambda}}

\newcommand{\sigg}{\sigma^2_{\bg}}
\newcommand{\Sigg}{\bSigma_{\bg}}
\newcommand{\Sigb}{\bSigma_{\bb}}
\newcommand{\sigx}{\sigma^2_{\bX}}
\newcommand{\sigy}{\sigma^2_{\bY}}
\newcommand{\siga}{\sigma^2_{\alpha}}
\newcommand{\sigb}{\sigma^2_{\bb}}
\newcommand{\sigbeta}{\sigma^2_{\bbeta}}

\newcommand{\Bqb}{B_{q(\sigma^2_{\bb})}}
\newcommand{\Bqg}{B_{q(\sigma^2_{\bg})}}
\newcommand{\Bqlk}{B_{q(\lambda_k)}}
\newcommand{\Bqx}{B_{q(\sigma^2_{\bX})}}
\newcommand{\Bqy}{B_{q(\sigma^2_{\bY})}}

\newcommand{\muqsigy}{\mu_{q(1/\sigma^2_{\bY})}}
\newcommand{\muqsigb}{\mu_{q(1/\sigma^2_{\bb})}}
\newcommand{\muqsigg}{\mu_{q(1/\sigma^2_{\bg})}}
\newcommand{\muqsigx}{\mu_{q(1/\sigma^2_{\bX})}}

\newcommand{\muqg}{\bmu_{q(\bg)}}
\newcommand{\Sigqg}{\bSigma_{q(\bg)}}

\newcommand{\muqc}{\bmu_{q(\bC)}}
\newcommand{\muqcij}{\bmu_{q(\bc), ij}}
\newcommand{\Sigqc}{\bSigma_{q(\bC)}}

\newcommand{\muqe}{\mu_{q(\varepsilon)}}
\newcommand{\Sigqe}{\Sigma_{q(\varepsilon)}}

\newcommand{\muqb}{\bmu_{q(\bb)}}
\newcommand{\Sigqb}{\bSigma_{q(\bb)}}

\newcommand{\muqbeta}{\bmu_{q(\bbeta)}}
\newcommand{\Sigqbeta}{\bSigma_{q(\bbeta)}}

\newcommand{\bbar}{\overline{\beta}}


\hypersetup{colorlinks,linkcolor=,urlcolor=MainColor}

\title{Lecture 10: GLMs: Poisson Regression, Overdispersion}
\author{Nick Reich / Transcribed by Daveed Goldenberg, edited by Josh Nugent}
\date{}

\begin{document}
\frame{\titlepage}

\begin{frame}{Poisson GLMs}

\begin{itemize}
\tightlist
\item
  Imagine you have count data following the Poisson distribution: \[
  Y_i \sim Poisson(\lambda_i)
  \]
\item
  \(Y_i\) is the total count in the time interval, \(\lambda_i\) is
  \(E(Y_i)\), that is, the risk/rate of occurrence in some time
  interval,\\
\item
  We use a log link for our GLM:
\end{itemize}

\[
\eta_i = X_i\beta=\log{\lambda_i}=g(\lambda_i)=g(E[y_i])
\]

\end{frame}

\begin{frame}{Poisson GLMs}

\begin{itemize}
\tightlist
\item
  Key Points:
\item
  Log link implies multiplicative effect of covariates
\end{itemize}

\[
\begin{aligned}
&log(\lambda) = \beta_0 + \beta_1X_1 + \beta_2X_2\\
& \lambda = e^{\beta_0}e^{\beta_1X_1}e^{\beta_2X_2}
\end{aligned}
\] - Relative risk is the interpretation for \(e^{\beta}\)

\[
\begin{aligned}
log(\lambda_i |X_1 = k+1, X_2 = c) = \beta_0 + \beta_1(k+1) + \beta_2(c) \\
- log(\lambda_i | X_1 = k, X_2 = c) = \beta_0 + \beta_1(k) + \beta_2(c) \\
log((\lambda_i |X_1 = k+1, X_2 = c) / \lambda_i |X_1 = k+1, X_2 = c) = \beta_1
\end{aligned}
\]

\end{frame}

\begin{frame}{Exposure / Offset Term}

Often Poisson models have an `exposure' or `offset' term, representing a
demoninator of some kind. Examples: Let \(u_i\) be offset for
\(Y_i\)\ldots{}

\begin{itemize}
\item
  Disease incidence: \(Y_i\) = the number of cases of flu in a
  population in 1 year (in location i), \(u_i\) = population size
\item
  Accident rates: \(Y_i\) = the number of traffic accidents at site i in
  1 day, \(u_i\) = average number of vehicles travelling through site i
  in 1 day, or \(u_i\) = the number of vehicles through site i yesterday
\item
  The offset is used to scale the \(Y_i\)
\end{itemize}

\end{frame}

\begin{frame}{Exposure / Offset Term}

\[
\begin{aligned}
Y_i \sim Poisson(u_i * \lambda_i) \\
E(Y_i) = u_i *  \lambda_i \\
log(E(Y_i)) = log(u_i) + log(\lambda_i)
\end{aligned}
\]

\begin{itemize}
\tightlist
\item
  \(log(u_i)\) is our offset (from observed data, can be thought of as
  an intercept)
\item
  log(\(\lambda_i\)) is our \(\eta_i\) (the linear predictor)
\end{itemize}

\end{frame}

\begin{frame}{Exposure / Offset Term}

\begin{itemize}
\tightlist
\item
  In R, the Poisson glm can be specified with an offset
\item
  glm(Y \(\sim X_1 + X_2\), family = `poison', offset = log(u), data
  \ldots{})
\item
  the log is important in order to get the correct offset
\item
  The offset term is adding more information to the model but not
  estimating a coefficient
\end{itemize}

\end{frame}

\begin{frame}{Exposure / Offset Term}

\[
\begin{aligned}
Y_i \sim Poisson(u_i * \lambda_i)
\end{aligned}
\]

\begin{itemize}
\tightlist
\item
  The \(Y_i\) could be cases per day
\item
  \(u_i\) could be population (persons)
\item
  then \(\lambda_i\) would be cases per day per population (persons)
\item
  which makes this a rate for an individual
\end{itemize}

\[
\begin{aligned}
log(E(Y_i)) - log(u_i) = log(\lambda_i) \\
log(E(Y_i)/ u_i) = log(\lambda_i)
\end{aligned}
\]

\end{frame}

\begin{frame}{Overdispersion}

\begin{itemize}
\item
  In Poisson models for \(Y_i \sim Poisson(\lambda_i)\) \[
  \begin{aligned}
  Var(Y_i) = \lambda_i
  \end{aligned}
  \]
\item
  In GLM estimation notation
\end{itemize}

\[
\begin{aligned}
\mu_i = E(\lambda_i) \\
Var(\mu_i) = \lambda_i
\end{aligned}
\]

\begin{itemize}
\tightlist
\item
  In an overdispersed model, the variance is higher because of some
  variability not captured by Poisson
\end{itemize}

\[
\begin{aligned}
Var(\mu_i) = \phi\lambda_i  \\ \phi > 0 
\end{aligned}
\] - Overdispersion implies \(\phi>1\)

\end{frame}

\begin{frame}{Overdispersion}

\begin{itemize}
\tightlist
\item
  Likelihood equations for Poisson GLM
\end{itemize}

\[
\begin{aligned}
\sum_{i=1}^{N} \frac{(y_i - u_i)x_{ij}}{Var(\mu_i)} \frac{\partial\mu_i}{\partial\eta_i} = 0 \\
j = 0, ..,p
\end{aligned}
\]

\begin{itemize}
\item
  Depends on the distribution of Y through \(\mu_i\) and Var(\(\mu\))
\item
  \(\phi\) drops out of the likelihood equations - this makes sense;
  variability won't affect the MLE - that is, \(\beta\)s are identical
  for models with \(\phi\) \textgreater{} 1 and \(\phi\) = 1
\item
  However, \(\phi\) does impact estimated standard errors
\end{itemize}

\end{frame}

\begin{frame}{Overdispersion}

\[
\begin{aligned}
w_i = (\frac{\partial u_i}{\partial \eta_i})^2 / Var(Y_i) \\
cov(\hat{\beta}) = (X^TWX)^{-1} = \phi cov(\hat{\beta})
\end{aligned}
\]

\begin{itemize}
\tightlist
\item
  \(\phi\) does not affect the \(\beta\)s but it does affect their
  covariance as a scaling factor
\end{itemize}

\end{frame}

\begin{frame}{Is Overdispersion Term Needed in a Model?}

\begin{itemize}
\item
  (See example 4.7.4 in Agresti)
\item
  Start with standardized residuals
\item
  Assume:\\
  \[
  \begin{aligned}
  z_i = \frac{y_i - \hat{y_i}}{\sqrt{Var(\hat{y_i})}} \\
  = \frac{y_i - \mu_i}{\sqrt{\mu_i}} \sim N(0,1) \\
  \sum_{i =1}^n z_i^2 \sim \chi^2_{n-k}
  \end{aligned}
  \]
\item
  where k is the number of parameters\\
\item
  if the sum of \(z_i^2\) is large (compare to chi-squared), we may need
  an overdispersion term \(\phi\)
\end{itemize}

\end{frame}

\begin{frame}{Is Overdispersion Term Needed in a Model?}

\[
\begin{aligned}
\hat{\phi} = \frac{\sum_{i=1}^n z_i^2}{n-k}
\end{aligned}
\]

\begin{itemize}
\item
  summarizes overdispersion in data compared to the fitted model
\item
  if \(\phi^2\) \textgreater{} 1, we should use the ``quasipoisson''
  family in R's glm() function
\item
  The SEs of a quasipoisson model are equivalent to the SEs of the
  Poisson model miultiplied by \(\sqrt{\hat{\phi}}\)
\end{itemize}

\[
\begin{aligned}
SE_{qp}(\hat{\beta}) = SE_p(\hat{\beta})*\sqrt{\hat{\phi}}
\end{aligned}
\]

\end{frame}

\end{document}
